\section{Methodology}

\subsection{Calibration Inversion Metric}

For each task, we compute calibration inversion score:

\begin{equation}
\text{inv}(t) = P(\text{wrong}_t) - P(\text{correct}_t)
\end{equation}

Tasks with $\text{inv}(t) > 0.1$ exhibit calibration inversion.

\subsection{Calibration Clustering}

We cluster tasks using K-means with silhouette validation. The optimal $k=2$ with silhouette$=0.6016$ reveals two behavioral clusters.

\subsection{Mechanism Verification Framework}

We decompose the mechanism into five testable sub-hypotheses:

\begin{table}[h]
\centering
\begin{tabular}{@{}lllll@{}}
\toprule
Hypothesis & Test & Threshold & Justification & Gate \\
\midrule
H-E1 & Silhouette score & $> 0.3$ & Standard clustering threshold \citep{rousseeuw1987silhouettes} & MUST\_WORK \\
H-M1 & Confidence overlap & $> 0.7$ or $< 0.1$ & Statistical comparison literature & MUST\_WORK \\
H-M2 & Rate difference & $< 0.15$ & Small effect size ($d < 0.2$) & SHOULD\_WORK \\
H-M3 & Separation score & $< 0.1$ & Low separation in representation space & SHOULD\_WORK \\
H-M4 & Correlation $r$ & $> 0.4$ & Moderate effect size (Cohen) & SHOULD\_WORK \\
\bottomrule
\end{tabular}
\caption{Mechanism verification framework with testable hypotheses.}
\label{tab:hypotheses}
\end{table}

\subsection{Bidirectional Feature Detection}

We use three keyword-based features: user-belief-reference, context-dependent, hedged-answer. Known limitation: achieved only 2.3\% prevalence.
